% ===== 1. Титульний =====
\begin{frame}[plain]
  \vspace*{-0.6cm}\titlepage
  \note{[Доповідач 1] Добрий день! Ми сьогодні захищаємо перший етап командного проєкту — вибір теми. Одразу скажемо результат: ми обрали тему «Цифровий тренер для автоспорту». Це вебзастосунок, який перетворює GPS-запис з екшн-камери чи телефона на розбір заїзду з порадами. Але до цього рішення ми дійшли не одразу. Ми розглянули дві ідеї, для кожної написали опис і зробили SWOT-аналіз, а потім порівняли їх за однаковими критеріями. Далі покажемо весь цей шлях по кроках. Говоритимемо по черзі, кожен про свою частину.}
\end{frame}

% ===== 2. Мета і процес =====
\begin{frame}{Мета етапу і як ми обирали тему}
  \textbf{Мета:} обрати тему командного проєкту й обґрунтувати вибір.

  \vspace{1.2em}
  \centering
  \begin{tikzpicture}[
      stp/.style={draw=accent,fill=accent!8,rounded corners=3pt,line width=0.8pt,
        minimum width=2.15cm,minimum height=1.1cm,align=center,font=\small\bfseries},
      final/.style={stp,fill=accent,text=white},
      desc/.style={font=\small,text width=2.3cm,align=center,text=ink!80},
      arr/.style={-{Stealth[length=2.2mm]},line width=0.9pt,draw=ink!60},
      node distance=0.5cm]
    \node[stp] (s1) {Ідеї};
    \node[stp,right=of s1] (s2) {Опис};
    \node[stp,right=of s2] (s3) {SWOT};
    \node[stp,right=of s3] (s4) {Порівняння\\за критеріями};
    \node[final,right=of s4] (s5) {Рішення};
    \foreach \a/\b in {s1/s2,s2/s3,s3/s4,s4/s5} \draw[arr] (\a) -- (\b);
    \node[desc,below=0.3cm of s1] {дві теми};
    \node[desc,below=0.3cm of s2] {що, для кого, яка проблема};
    \node[desc,below=0.3cm of s3] {сильні й слабкі сторони, можливості, загрози};
    \node[desc,below=0.3cm of s4] {5 критеріїв з вагами, оцінки 1–5};
    \node[desc,below=0.3cm of s5] {цифровий тренер для автоспорту};
  \end{tikzpicture}

  \note{[Доповідач 1] Мета етапу — не просто назвати тему, а обґрунтувати, чому саме її. Тому ми йшли за простою схемою з п'яти кроків. Спочатку зібрали ідеї: у підсумку їх дві — цифровий тренер для автоспорту і розумний кошик для товарів. Для кожної написали опис: що це, для кого і яку проблему вирішує. Потім зробили SWOT-аналіз із джерелами. Далі домовилися про п'ять критеріїв з вагами і поставили кожній темі оцінки від одного до п'яти. І лише після цього ухвалили рішення. Так обидві теми ми порівнювали за однаковими правилами, а не за враженням.}
\end{frame}
